\section*{अध्याय \ref{ch:b2:global-change} --- वैश्विक परिवर्तन और जीवमंडल}

\begin{solution}{exo:b2:global-change:1}
$\Delta T = \lambda C_{\text{संचयी}}$, जिसमें $\lambda = \qty{1.65}{\celsius}$
प्रति \qty{1000}{GtC}: \qty{1.2}{\celsius}, \qty{1.65}{\celsius},
\qty{3.3}{\celsius}।
\end{solution}

\begin{solution}{exo:b2:global-change:2}
अंश-दिन: दैनिक माध्य ताप की किसी आधार से अधिकता का
दिनों भर योग, यानी वह \omterm{thm:b2:global-change:phenology}{तापीय समय} जो पौधों तथा बहिस्तापियों का
विकास बाँधता है। बेमेल: जब दो अन्योन्यक्रिया करती जातियाँ अपनी ऋतुएँ
भिन्न संकेतों से समयबद्ध करें और तापन एक को चला दे और दूसरी को नहीं ---
यानी वह मक्खीमार, जो अफ़्रीका में दिन-लंबाई से संकेतित है, अपनी पुरानी
तिथि पर पहुँचकर उस इल्ली-शिखर को, जो ताप से संकेतित है, दो सप्ताह बीता हुआ पाता है।
\end{solution}

\begin{solution}{exo:b2:global-change:3}
प्रति अंश लगभग \qty{150}{m} ऊपर और \qty{150}{km} ध्रुव की ओर। शिखर पर
की किसी जाति के पास कोई ऊँची ज़मीन नहीं: उसकी समताप-रेखा उस पर्वत से
ऊपर चढ़ जाती है और उसका आवास मिट जाता है, जबकि ऊपर जाते-जाते हर
पट्टी का क्षेत्रफल ऊँचाई के साथ सिकुड़ता है।
\end{solution}

\begin{solution}{exo:b2:global-change:4}
घुली \ensuremath{\mathrm{CO_2}} कार्बोनिक अम्ल बनाती है, जो
हाइड्रोजन आयन छोड़ता है, और वह pH \ensuremath{\mathrm{CO_2}} दाब के
लघुगणक के समानुपात में गिरता है। वे अतिरिक्त
हाइड्रोजन आयन कार्बोनेट को बाइकार्बोनेट में बदल देते हैं, इसलिए
कार्बोनेट-सांद्रता तथा कैल्सियम कार्बोनेट की \omterm{thm:b2:global-change:acid}{संतृप्ति अवस्था} गिरती है।
जो जीव ऐरागोनाइट बनाते हैं --- प्रवाल, टेरोपॉड, बहुत-से डिंभक
--- और कैल्साइट --- कोकोलिथोफ़ोर, फ़ोरामिनिफ़ेरा, मोलस्क ---
उन्हें कवच बनाने में अधिक ख़र्च करना पड़ता है और, संतृप्ति से नीचे, वे उन्हें खो देते हैं।
\end{solution}

\begin{solution}{exo:b2:global-change:5}
\qty{1.5}{\celsius}: कुल मिलाकर $1.5/1.65\times 1000 = \qty{909}{GtC}$,
यानी \qty{209}{GtC} शेष, यानी 21 वर्ष। \qty{2}{\celsius}:
\qty{1212}{GtC}, \qty{512}{GtC} शेष, यानी 51 वर्ष।
\end{solution}

\begin{solution}{exo:b2:global-change:6}
$\sqrt{2\times 150/0.12} = 50$ दिन; और जल्दी आना $1.5/0.12 = 12.5$
दिन।
\end{solution}

\begin{solution}{exo:b2:global-change:7}
\qty{3}{\celsius}: समताप-रेखाएँ $3000/6.5 = \qty{460}{m}$ चढ़ती हैं; वह
पट्टी खिसककर \qtyrange{1260}{1860}{m} पर आती है, पर वह पर्वत
\qty{1800}{m} पर ख़त्म हो जाता है: उस पट्टी के ऊपर के \qty{60}{m} खो जाते हैं और बाक़ी
वहाँ बैठती है जहाँ क्षेत्रफल पहले का $2^{-460/300} = 0.35$ है --- यानी
लगभग दो तिहाई क्षेत्रफल गया। \qty{5}{\celsius}: \qty{770}{m}; वह
पट्टी \qtyrange{1570}{2170}{m} होती; पर केवल \qtyrange{1570}{1800}{m}
है, यानी उस पट्टी की चौड़ाई का एक तिहाई, और वह भी क्षेत्रफल के
$2^{-770/300} = 0.17$ पर --- यानी मूल क्षेत्रफल के दसवें भाग से भी कम: वह जाति
प्रभाव में जा चुकी है।
\end{solution}

\begin{solution}{exo:b2:global-change:8}
560: pH $7.93$, हाइड्रोजन आयन $\times 1.9$; 800: $7.79$, $\times
2.6$; 1000: $7.70$, $\times 3.1$।
\end{solution}

\begin{solution}{exo:b2:global-change:9}
$1 - (1 - h)^{0.25}$: \qty{26}{\%}, \qty{53}{\%}, \qty{82}{\%}। वह
तुरंत हानि इसलिए छोटी है कि वह अवशेष अब भी अधिकांश जातियों की
समष्टियाँ थामे है; वे टिकने को बहुत छोटी हैं और
दशकों भर मिट जाती हैं --- यानी वह विलोपन-ऋण।
\end{solution}

\begin{solution}{exo:b2:global-change:10}
$c = 2\sqrt{0.7\times 20} = \qty{7.5}{km/yr}$; \qty{1000}{km}
लगभग 130 वर्षों में। $r$ या $D$ में से किसी को भी आधा करना $c$ को
$\sqrt 2$ से बाँट देता है, यानी \qty{5.3}{km/yr} तक; कोई भी उसे आधा नहीं करता --- उस
गति को आधा करने के लिए गुणनफल $rD$ को एक चौथाई करना पड़ेगा।
\end{solution}

\begin{solution}{exo:b2:global-change:11}
पृष्ठभूमि: प्रति वर्ष $8\times 10^{6}\times 10^{-6} = 8$, प्रति सदी 800;
और हज़ार गुना पर, प्रति वर्ष 8000, प्रति सदी \num{800000}।
सदी भर में दस हज़ार प्रलेखित लगभग पृष्ठभूमि से दस गुना है,
हज़ार नहीं --- पर प्रलेखन उन जातियों के केवल
कुछ प्रतिशत को छूता है जो वर्णित तथा निगरानी में हैं (पक्षी,
स्तनी, कुछ पौधे), जिनमें वह दर सचमुच पृष्ठभूमि से सौ
गुना या अधिक है; और उस अवर्णित बहुसंख्या के लिए अधिकांश
विलोपन कभी देखे ही नहीं जाते, और वह जाति--क्षेत्रफल आकलन ही
एकमात्र पकड़ है। सच गिने हुए तथा
आँके हुए के बीच है, और इसीलिए बताई गई परास इतनी चौड़ी है।
\end{solution}

\begin{solution}{exo:b2:global-change:12}
वह ताप \omterm{prop:b2:global-change:tcre}{संचयी उत्सर्जनों} के समानुपाती है; और जब
शुद्ध उत्सर्जन शून्य हों तब वह संचयी योग बढ़ना बंद कर देता है और वह
तापन भी। उस \omterm{thm:b2:biogeochemical-cycles:box}{डिब्बा-प्रतिरूप} में, उत्सर्जन शून्य होने पर, वह
वायुमंडलीय अतिरेक गिरने लगता है ज्यों-ज्यों वे पात्र काम करते रहें, जो
ठंडा करता --- पर वह महासागर, जो चढ़ी हुई \ensuremath{\mathrm{CO_2}} के साथ
साम्य की ओर अब भी गरम हो रहा है, उस सतह को गरम करता रहता है,
और वे दोनों प्रभाव लगभग कट जाते हैं: वह ताप शताब्दियों तक जहाँ है
वहीं रहता है। उस तर्क में से ग़ायब: बाक़ी
हरितगृह गैसें (मीथेन का छोटा जीवनकाल कहता है कि उत्सर्जन रुकते ही उसका तापन
जल्दी गिरता है), वे एरोसॉल जो अभी कुछ तापन को ढकते हैं और
उन्हीं उत्सर्जनों के सप्ताहों के भीतर मिट जाते हैं जो उन्हें बनाते
हैं, और वे प्रतिपुष्टियाँ --- स्थायी-तुषार का कार्बन, वन-मरण ---
जो अपने ही उत्सर्जन जोड़ सकती हैं।
\end{solution}

\begin{solution}{pb:b2:global-change:1}
\textbf{1.} A: $700 + 80\times 10 = \qty{1500}{GtC}$। B: $700 +
\tfrac12\times 50\times 10 = \qty{950}{GtC}$।
\textbf{2.} A: $1.65\times 1.5 = \qty{2.5}{\celsius}$। B:
$1.65\times 0.95 = \qty{1.6}{\celsius}$।
\textbf{3.} 2050. A: \qty{1000}{GtC}, \qty{1.65}{\celsius}। B:
उत्सर्जन $10\int_0^{30}(1 - t/50)\,\mathrm{d}t = 10(30 - 9) =
\qty{210}{GtC}$, कुल 910, \qty{1.5}{\celsius}।
\textbf{4.} प्रति दशक $\lambda\times 100 = \qty{0.165}{\celsius}$,
यानी देखे गए $0.2$ से थोड़ा नीचे (जिसमें बाक़ी गैसें
और एरोसॉल-आवरण का खोना भी शामिल हैं)।
\textbf{5.} हाँ: \omterm{prop:b2:global-change:tcre}{संचयी उत्सर्जन} 2070 में बढ़ना बंद कर देते हैं, और वह
तापन भी, जो \qty{1.6}{\celsius} पर रह जाता है --- यानी वह
समानुपातिता कहती है कि वह ताप उस कुल से बैठता है, उस
दर से नहीं।
\textbf{6.} $2/1.65\times 1000 = \qty{1212}{GtC}$; और पथ A उस तक
$512/10 = 51$ वर्ष बाद, यानी 2071 में पहुँचता है।
\textbf{7.} $\sqrt{2\times 200/0.1} = 63$ दिन।
\textbf{8.} 2020 के सापेक्ष तापन: A $2.5 - 1.2 =
\qty{1.3}{\celsius}$, B \qty{0.4}{\celsius}। जल्दी आना $\Delta T/b$:
A 13 दिन, B 4 दिन।
\textbf{9.} अंतर: A $20 + 13 = 33$ दिन; B 24 दिन।
\textbf{10.} इल्लियाँ पत्ती निकलने के बाद दिन 15 से दिन 25 तक
उपलब्ध हैं। B: वह पक्षी दिन 24 पर पहुँचता है, यानी उस खिड़की में, बस।
A: दिन 33, यानी अंतिम इल्ली के आठ दिन बाद।
\textbf{11.} वह पक्षी $3\times 1.3 = 4$ दिन जल्दी आता है: अंतर $33 - 4 =
29$ दिन --- यानी अब भी उस खिड़की के बाहर।
\textbf{12.} वह पक्षी नए ताप पर जी सकता है; जो वह नहीं
कर सकता वह है भोजन के जा चुकने पर चूज़े पालना, क्योंकि उसका संकेत
और उसके शिकार का संकेत तापन का उत्तर भिन्न ढंग से देते हैं।
\textbf{13.} A: $1.3/6.5\times 1000 = \qty{200}{m}$ ऊपर और
$1.3/0.65\times 100 = \qty{200}{km}$ ध्रुव की ओर। B: \qty{57}{m} तथा
\qty{57}{km}।
\textbf{14.} 80 वर्षों में वह वृक्ष \qty{8}{km} चलता है: वह
\qty{200}{km} (A) के पीछे नहीं चल सकता और \qty{57}{km} (B) से भी कम पड़ता है ---
यानी वृक्ष पिछड़ते हैं, और उनके परास दोनों पथों पर शताब्दियों तक
जलवायु के साथ साम्य से बाहर रहेंगे।
\textbf{15.} A: वह पट्टी \qtyrange{1700}{2000}{m} होती: वह
बस शिखर तक पहुँचती है, और आगे जाने को कहीं नहीं। वह पर्वत ऊँचाई के साथ
सँकरा होता है, इसलिए वही टुकड़ा अपने पुराने क्षेत्रफल का केवल $2^{-200/300} = 0.63$ गुना
ढकता है --- और किसी भी आगे के तापन के पास उसे चलाने को कोई ज़मीन नहीं बची।
\textbf{16.} B: \qty{57}{m} ऊपर, क्षेत्रफल-गुणक $2^{-57/300} = 0.88$: यानी
\qty{12}{\%} की हानि।
\textbf{17.} \qty{30}{km} पर आठ दशक: \qty{240}{km}, यानी A पर ज़रूरी
\qty{200}{km} से अधिक: वह मछली साथ चलती रहती है (और अपनी पुरानी
मछली-भूमि से बाहर चली जाती है)।
\textbf{18.} निचले किनारे पर वह जाति ऐसे होड़ियों,
रोगजनकों तथा परभक्षियों से मिलती है जो नीचे से ऊपर चल रहे हैं और जिनका सामना उसने पहले नहीं किया,
और उसकी अपनी शरीरक्रिया, जो पुराने ताप के अनुरूप ढली है,
उनकी से हार जाती है: यानी वह पीछे हटना जितना शरीरक्रियात्मक है उतना ही होड़ का।
\textbf{19.} A: $\mathrm{pH} = 8.2 - 0.9\log_{10}(600/280) = 7.90$;
हाइड्रोजन आयन 1750 के $10^{0.30} = 2.0$ गुना। B: $8.04$; $1.44$
गुना।
\textbf{20.} $\Omega = 4/(\text{अनुपात})$: A $2.0$, B $2.8$। 2 पर वे
प्रवाल अब भी कैल्सीभवन करते हैं पर धीरे, उनके कंकाल कमज़ोर और उनकी
भित्तियाँ बढ़ने से तेज़ अपरदित; और $2.8$ पर वे बनाते हैं, यद्यपि
पथ B की ऊष्मा फिर भी उन्हें विवर्णित करती है।
\textbf{21.} A: $400\times 0.4^{0.25} = 318$, यानी 82 खोईं। B:
$400\times 0.8^{0.25} = 378$, यानी 22 खोईं।
\textbf{22.} A: $318\times(1 - 0.05\times 1.3) = 298$। B:
$378\times(1 - 0.05\times 0.4) = 371$।
\textbf{23.} A: $1 - 298/400 = \qty{26}{\%}$। B: \qty{7}{\%}।
\textbf{24.} दोनों पर, सफ़ाई: A पर 20 के मुक़ाबले 82 जातियाँ, B पर 7
के मुक़ाबले 22 --- यानी उस प्रतिरूप का तापन-पद संयत है; असल में
वे दोनों अन्योन्यक्रिया करते हैं, क्योंकि कोई खंडित वन अपना परास खिसका नहीं सकता।
\textbf{25.} पथ A: \qty{2.5}{\celsius}, pH $7.90$, यानी वृक्ष-जातियों का
एक चौथाई खोया। पथ B: \qty{1.6}{\celsius}, pH $8.04$, \qty{7}{\%}
खोया।
\end{solution}
